\scanpage{123}
\mspage{61}
in these parameters, and reducing to a relatively prime $f,g$, \unsure{one is|one sees} --- blowing up $a$ once more --- that the divisors of $f,g$ \deleted{\unsure{are|themselves}} \unsure{separate}\ldots

b) A valuation ring $V$ \deleted{which} for the extension $L/K$ dominates a well-determined point of each $X$, whence an $x\in\mathfrak X=\varprojlim X$, and $V$ dominates $\underline{O}_{\mathfrak X,x}$, hence is equal to it since the latter is a valuation ring.

\uline{Note} Every \inserted{model $U$ of $L/K$ defines an} \uline{open} of $\mathfrak X$ \deleted{\unsure{corresponding}} \deleted{\unsure{obviously}} \inserted{(by embedding it} in a proper model $X$, this is clear) and \deleted{and} the models correspond to \uline{certain} opens of $\mathfrak X$ (\uline{not} stable under finite unions).

For every \inserted{\unsure{norma\textsuperscript{l}|normal}} \unsure{loca\textsuperscript{l}|local} ring $O\subset L$ with field of fractions $L$, one can look at $\operatorname{Spec}\underline{O}=S$, its closed point $\unsure{x|s}$, the \unsure{normali-}
\margin{122}
